\chapter*{KATA PENGANTAR} \addcontentsline{toc}{chapter}{KATA PENGANTAR}
\begin{singlespace}
Puji dan syukur penulis panjatkan ke hadirat Tuhan Yang Maha Esa 
atas segala rahmat, berkat, dan karunia-Nya, sehingga Laporan Kemajuan 
Magang Industri ini dapat diselesaikan tepat pada waktunya. Laporan ini disusun sebagai salah satu syarat kelengkapan administrasi 
dan akademik dalam program Merdeka Belajar Kampus Merdeka (MBKM) 
Magang Industri bagi mahasiswa semester enam di Program Studi Teknik 
Informatika, Fakultas Teknologi Informasi, Universitas Tarumanagara. 

Pelaksanaan magang yang dilakukan di PT Inti Utama Solusindo 
(bagian dari Pharos Group) merupakan upaya strategis untuk menjembatani 
kompetensi teoritis yang diperoleh di bangku perkuliahan dengan 
kebutuhan praktis di industri teknologi informasi yang dinamis.

Penulis menyadari bahwa keberhasilan pelaksanaan magang dan 
penyusunan laporan ini tidak terlepas dari bimbingan, arahan, 
dan dukungan berbagai pihak. Oleh karena itu, penulis ingin 
menyampaikan apresiasi dan terima kasih yang sebesar-besarnya kepada:

\begin{enumerate}[label=\arabic*.,
    leftmargin=2.5em,
    labelindent=0pt,]
    \item Ibu Prof. Lina, S.T., M.Kom., Ph.D., selaku Dekan Fakultas Teknologi 
          Informasi Universitas Tarumanagara, yang telah memberikan dukungan 
          terhadap pelaksanaan program Merdeka Belajar Kampus Merdeka di 
          lingkungan Fakultas Teknologi Informasi.
    \item Ibu Lely Hiryanto, S.Kom., M.Kom., Ph.D., selaku Ketua Program Studi 
          Teknik Informatika Fakultas Teknologi Informasi Universitas 
          Tarumanagara, yang telah memberikan arahan akademik bagi mahasiswa 
          dalam pelaksanaan kegiatan magang industri.
    \item Bapak Tony, S.Kom., M.Kom., Ph.D., selaku Koordinator Magang Industri 
          sekaligus Manajer Kerja Sama dan Sumber Daya Fakultas Teknologi 
          Informasi Universitas Tarumanagara, yang telah memfasilitasi dan 
          mendukung pelaksanaan program magang industri bagi mahasiswa.
    \item Ibu Viny Christanti Mawardi, S.Kom., M.Kom., selaku dosen pembimbing 
          yang telah memberikan bimbingan, arahan, serta masukan selama 
          pelaksanaan magang dan penyusunan laporan ini.
    \item Saudara Dimas Aditya Nugroho, selaku mentor di PT Inti Utama Solusindo, 
          yang telah memberikan arahan, pengetahuan, serta pengalaman praktis 
          selama proses magang berlangsung.
    \item Rekan-rekan kerja di PT Inti Utama Solusindo dari Fakultas Teknologi 
          Informasi Universitas Tarumanagara yang telah memberikan dukungan, 
          kerja sama, serta berbagi pengalaman selama pelaksanaan kegiatan 
          magang.
\end{enumerate}

Penulis menyadari bahwa laporan kemajuan ini masih jauh dari kesempurnaan, 
baik dari segi kedalaman analisis maupun teknis penulisan. Terdapat 
berbagai tantangan yang masih harus dihadapi dalam sisa periode magang,
terutama dalam hal penguasaan kerangka kerja yang lebih kompleks. 
Namun, penulis berkomitmen untuk terus belajar dan memberikan kontribusi
terbaik bagi perusahaan. Semoga laporan ini dapat memberikan manfaat 
bagi para pembaca, khususnya bagi rekan-rekan mahasiswa yang akan menempuh 
program serupa di masa mendatang.

\begin{flushright}
    Jakarta, \textcolor{red}{xx April 2026}\\
    Peserta Magang Industri,\\[2cm]
    Aldian Yohanes\\
\end{flushright}
\end{singlespace}
